\begin{example}
On reprend l'exemple \ref{exsem2} et on le reecrit comme un produit des matrices
$$
\gcd(f_{0},f_{1})=\begin{pmatrix}
0 & 1 \\
1 & -q_{3}
\end{pmatrix}\begin{pmatrix}
0 & 1 \\
1 & -q_{2}
\end{pmatrix}\begin{pmatrix}
0 & 1 \\
1 & -q_{1}
\end{pmatrix}\begin{pmatrix}
f_{0} \\
f_{1}
\end{pmatrix}
$$
donc
$$
(x^{2}+2)=\gcd(f_{0},f_{1})=(x+1)f_{0}+(2x^{3}+3x^{2}+4x)f_{1}
$$
qui est la premiere composante du produit matriciel.
\end{example}

\section{Factorisation en irreductibles}
\begin{definition}[Polynome irreductible]
Un polynome $p(x)\in K[x]$ est irreductible si :
\begin{enumerate}
\item $\deg p \geq 1$.
\item Si pour $f,g\in K[x]$ on a que $fg=p$, alors $\deg f = \deg p$ ou $\deg g =\deg p$.
\end{enumerate}
\end{definition}
\begin{example}
Tout polynome $ax+b\in K[x]$ tel que $a\neq 0$ est irreductible car si $ax+b=fg$ on a que
$$
\deg (ax+b)=1=\deg f+\deg g
$$
alors $f\in K$ ou $g\in K$.
\end{example}
\begin{example}
Le polynome $x^{2}+x+1\in \mathbf{Z}_{2}[x]$ est irreductible. Sinon ils existent $a,b\in \mathbf{Z}_{2}$ tels que
$$
x^{2}+x+1=(x-a)(x-b)
$$
mais ce polynome n'a pas de racines dans $\mathbf{Z}_{2}$ donc cet factorisation n'est pas possible.
\end{example}
\begin{theorem}
Soit $f(x)\in K[x]$ de $\deg f=\{ 2,3 \}$ alors $f$ est irreductible si et seulement si $f$ ne possede pas de racines dans $K$. 
\end{theorem}
\begin{proof}
\begin{itemize}
\item Sens $\Rightarrow$: Soit $\alpha \in K$ une racine de $f$ alors $(x-a)\mid f$ qui pose une contradiction avec l'irreductibilite.
\item Sens $\Leftarrow$: Comme $\deg f=\{ 2,3 \}$ si ils existent $a,b\in K[x]$ tels que $f=ab$ alors $\deg a+\deg b=\{ 2,3 \}$. Sans perte de generalite on pose que $\deg a=1$ et $a(x)=\beta(x-\alpha)$ avec $\alpha,\beta \in K$ alors $f(\alpha)=a(\alpha)b(\alpha)=0b(\alpha)=0$ qui fera de $\alpha$ une racine de $f$ posant une contradiction.
\end{itemize}
\end{proof}

Soit $f(x)\in K[x]$ et $\deg f \ge 1$ si $f$ n'est pas irreductible alors $f=p_{1}p_{2}$ avec $\deg p_{1},\deg p_{2} \ge 1$ et $\deg p_{1}+\deg p_{2}=\deg f$. On peut continuer a factoriser les polynomes $p_{i}$ pour obtenir:
\begin{equation}
f=a\prod_{i=1}^{\ell} p_{i} \label{decomp}
\end{equation}
ou $p_{i}\in K[x]$ sont irreductibles et unitaires et $a\in K$. Cette factorisation est unique.

\begin{lemma}\label{lem1}
Soit $p \in K[x]$ irreductible et $f_{1},\dots,f_{\ell}\in K[x]$. Si 
$$
p\mid f_{1}f_{2}\dots f_{\ell}
$$
il existe $i\in \{ 1,\dots,\ell \}$ tels que $p\mid f_{i}$.
\end{lemma}
\begin{proof}
Il suffit de montrer l'assertion pour $\ell=2$ car soit $p\mid f_{1}$ soit $p\mid f_{2}\dots f_{\ell}$. 

Soit $\ell=2$ et $p\mid f_{1}f_{2}$. Supposons que $p\nmid f_{1}$ alors $\gcd(p,f_{1})=1$ et ils existent $u,v\in K[x]$ tels que $1=up+vf_{1}$. Si on multiplie de deux cotes par $f_{2}$ on obtient
$$
f_{2}=upf_{2}+vf_{1}f_{2}
$$
et comme $p\mid f_{1}f_{2}$ alors il existe $g\in K[x]$ tel que $pg=f_{1}f_{2}$ :
$$
f_{2}=p(uf_{2}+vg)
$$
et $p\mid f_{2}$.
\end{proof}
\begin{theorem}
La decomposition \ref{decomp} est unique a l'ordre des $p_{i}$ pres.
\end{theorem}
\begin{proof}
Soient
$$
a\prod_{i=1}^{\ell} p_{i}=b\prod_{j=1}^{k} q_{j}
$$
deux factorisations possibles de $f$. Le fait que $p_{i}$ et $q_{j}$ sont unitaires implique que $a=b$. On peut utiliser \ref{lem1} pour avoir que $p_1\mid q_{j}$ pour un $j\in \{ 1,\dots ,k \}$ mais comme $p_{i}$ et $q_{j}$ sont unitaires et irreductibles alors $p_{i}=q_{j}$.
$$
\prod_{i=2}^{\ell} p_{i}=\prod_{\mu=1,\mu\neq j}^{k} q_{\mu}
$$
et on conclut par recurrence.
\end{proof}

\section{Le polynome minimal de \texorpdfstring{$A\in K^{n\times n}$}{A}}

\begin{definition}
Soit $A\in K^{n\times n}$ et $p(x)\in K[x]$. L'evaluation de $p(x)$ en $A$ est la matrice 
$$
p(A)=a_{0}\mathrm{Id}+a_{1}A+a_{2}A^{2}+\dots+a_{n}A^{n}\in K^{n\times n}
$$
\end{definition}
\begin{exercise}
Pour tous $f,g\in K[x]$ on a :
\begin{enumerate}
\item $(f+g)(A)=f(A)+g(A)$
\item $(fg)(A)=f(A)g(A)$
\item $1(A)=\mathrm{Id}_{n}$
\end{enumerate}
En autres mots l'application $\phi:K[x]\to K^{n\times n},p(x)\mapsto p(A)$ est un morphisme d'anneaux. 
\end{exercise}
\begin{proof}[Solution]
\begin{enumerate}
\item On a que 
$$
f(x)=a_{0}+a_{1}x+\dots +a_{n}x^{n}
$$
et
$$
g(x)=b_{0}+\dots b_{n}x^{n}
$$
donc 
$$
(f+g)(x)=(a_{0}+b_{0})+\dots +(a_{n}+b_{n})x^{n}
$$
alors

\begin{align*}
(f+g)(A) & =(a_{0}+b_{0})\mathrm{Id}_{n}+\dots +(a_{n}+b_{n})A^{n} \\
 & =a_{0}\mathrm{Id}_{n}+\dots +a_{n}A^{n}+b_{0}\mathrm{Id}_{n} +\dots +b_{n}A^{n} \\
 & =f(A)+g(A)
\end{align*}

\end{enumerate}
\end{proof}

\begin{definition}
Un polynome $p(x)\in K[x]$ annule $A\in K^{n\times n}$ si $p(A)=0_{n}\in K^{n\times n}$.
\end{definition}
\begin{theorem}\label{thm2}
Soit $p(x)\in K[x]\setminus \{ 0 \}$ un polynome de degre minimal parmi les polynomes qui annulent $A\in K^{n\times n}$. Alors $p(x)\mid h(x)$ pour tout $h(x)\in K[x]$ qui annule $A$. $p(x)$ est unique a une multiplication par un scalaire pres. Si $p(x)$ est unitaire alors $p(x)$ est unique.  
\end{theorem}
\begin{proof}
Soit $h(x)\in K[x]$ tel que $h(A)=0$. On divise $h(x)$ avec rest pour obtenir $h(x)=q(x)p(x)+r(x)$ avec $\deg r<\deg p$. Si on evalue en $A$ on obtient $h(A)=q(A)p(A)+r(A)=0$ donc $r(A)=0$ et $p(x)\mid h(x)$.
\end{proof}
\begin{exercise}[Exo etoile serie 2]
Soit $K$ un corps et $A\in K^{n\times n}$ une matrice inversible. Posons $\tilde{A}:=(\mathrm{com}A)^{T}$ la transposee de la comatrice de $A$.
\begin{enumerate}
\item Montrer que $\det(A)\neq 0$.
\item Montrer que $A^{-1}=\det(A)^{-1}\tilde{A}$.
\item Montrer que pour toute matrice $A\in \mathbf{Z}^{n\times n}$, il existe $B\in \mathbf{Z}^{n\times n}$ telle que $AB=I$ si et seulement si $\det(A)=\pm1$.
\end{enumerate}
\begin{reminder}
$\mathrm{com}(A)_{ij}=(-1)^{i+j}\det(A_{ij})$ ou $A_{ij}$ est la matrice dont on a enlevee la $i$-eme ligne et $j$-eme colonne.
\end{reminder}
\end{exercise}
\begin{proof}[Solution]
\begin{enumerate}
\item Comme $A$ est inversible on a que $AA^{-1}=I_{n}$ alors $\det(AA^{-1})=1$ donc par proprietes du determinant on a $\det(A)\det (A^{-1})=1$ et $\det(A)\neq 0$.
\item Pour conclure, il suffit de montrer que $A\tilde{A}=\det(A)I_{n}$ (par unicite de l'inverse d'une matrice)( Pour etre precis, il faut aussi demontrer que $\tilde{A}A=\det(A)I_{n}$). Considerons les coefficients diagonaux de $A\tilde{A}$, pour tout $i\in \{ 1,\dots ,n \}$ on a 

\begin{align*}
(A\tilde{A})_{ii} & =\sum_{j=1}^{n} a_{ij}\tilde{a}_{ji} \\
 & =\sum_{j=1}^{n} a_{ij}(\mathrm{com}A)_{ij} \\
 & =\sum_{j=1}^{n} a_{ij}(-1)^{i+j}\det(A_{ij}) \\
 & =\det A.
\end{align*}

Considerons les coefficients hors diagonal. Pour $i\neq j\in \{ 1,\dots,n \}$ on a

\begin{align*}
(A\tilde{A})_{ij} & =\sum_{k=1}^{n} a_{ik}\tilde{a}_{kj} \\
 & =\sum_{k=1}^{n} a_{ik}(\mathrm{com}A)_{jk} \\
 & =\sum_{k=1}^{n} a_{ik}(-1)^{j+k}\det(A_{jk}) \\
 & =\text{le determinant de la matrice }A\text{ dont on remplace la }\\&j\text{-eme ligne par une copie de la }i\text{-eme ligne} \\
 & =0
\end{align*}
car $\det$ est une forme alternee.
\item Pour montrer le sens $\Rightarrow$, on suppose qu'il existe $B\in \mathbf{Z}^{n\times n}$ tel que $AB=I_{n}$ donc $\det(AB)=1$ et comme dans $\mathbf{Z}$ les seuls diviseurs de 1 sont 1 et $-1$ alors on observe que $\det(A)=\pm1$.

On montre l'autre sens $\Leftarrow$, on suppose que $\det(A)=\pm1$. Par les memes techniques que dans 2 on voit que $A\tilde{A}=\det(A)I_{n}$ avec $\tilde{A}=(\mathrm{com}A)^{T}\in \mathbf{Z}^{n\times n}$. Comme $\det(A)=\pm 1$ on conclut que $B=\det(A)^{-1}\tilde{A}\in \mathbf{Z}^{n\times n}$.
\end{enumerate}
\end{proof}
\begin{exercise}[Exo plus serie 3]
Soit $V$ un espace vectoriel reel et soit $B=\{ v_{1},\dots,v_{4} \}$ une base de $V$.
\begin{enumerate}
\item Soit $f$ l'endomorphsime defini par
$$
f(v_{1})=v_{1}-v_{2},~f(v_{2})=2v_{2}-6v_{3},~f(v_{3})=-2v_{1}+2v_{2},~f(v_{4})=v_{2}-3v_{3}+v_{4}
$$
Ecrivez la matrice $A_{B}$ de l'application $f$ dans la base $B$. Est-ce que $f$ est inversible? Si oui, ecrivez la matrice $A^{-1}_{B}$ de l'application inverse $f^{-1}:V\to V$.
\end{enumerate}
\end{exercise}
\begin{proof}[Solution]
On construit la matrice $A_{B}$:
$$
A_{B}=\begin{pmatrix}
1 & 0 & -2 & 0 \\
-1 & 2 & 2 & 1 \\
0 & -6 & 0 & -3 \\
0 & 0 & 0 & 1
\end{pmatrix}
$$
on peut observer que il y a une relation de dependence entre la troisieme et premiere colonne donc $A_{B}$ n'est pas inversible.
\end{proof}

